% Part: first-order-logic
% Chapter: axiomatic-proofs
% Section: identity

\documentclass[../../../include/open-logic-section]{subfiles}

\begin{document}

\olfileid{fol}{axd}{ide}

\olsection{\usetoken{P}{derivation} نال \usetoken{S}{identity}}

!!{derivation}s وچ $\eq$ نوں شامل کرن لئی اسیں بس نویں مسلّمے دے
خاکے شامل کردے آں۔ !!{derivation} تے $\Proves$ دیاں تعریفاں
اوہو رہندیاں نیں؛ ہُن نویں مسلّماں نوں وی اجازت اے۔

\begin{defn}[!!{identity} دے مسلّمے]
\begin{align}
\ollabel{ax:id1} & \eq[t][t], \\
\ollabel{ax:id2} & \eq[t_1][t_2] \lif (!B(t_1) \lif
  !B(t_2)),
\end{align}
% PNBAXD-010: ماخذ دے ’’بند ٹرم‘‘ نوں اگلیاں بے قید قضیاں نال ملا کے ایتھے عام ٹرم کیتا۔
کسے وی ٹرماں $t$, $t_1$, $t_2$ لئی۔
\end{defn}

\begin{prop}\ollabel{prop:sound}
  مسلّمے \olref{ax:id1} تے \olref{ax:id2} منطقی طور تے معتبر نیں۔
\end{prop}

\begin{proof}
  مشق۔
\end{proof}

\begin{prob}
\olref[fol][axd][ide]{prop:sound} ثابت کرو۔
\end{prob}

\begin{prop}\ollabel{prop:iden1}
 $\Gamma \Proves \eq[t][t]$، کسے وی ٹرم $t$ تے سیٹ~$\Gamma$ لئی۔
\end{prop}

\begin{prop}\ollabel{prop:iden2}
  جے $\Gamma \Proves !A(t_1)$ تے $\Gamma \Proves
  \eq[t_1][t_2]$، تاں $\Gamma \Proves !A(t_2)$۔
\end{prop}

\begin{proof}
!!{formula}
\[
(\eq[t_1][t_2] \lif (!A(t_1) \lif !A(t_2)))
\]
~\olref{ax:id2} دی اک مثال اے۔ نتیجہ \MP نوں دو واری
لاون نال نکلدا اے۔
\end{proof}

\end{document}
